\documentclass[11pt]{article}
\usepackage[utf8]{inputenc}
\usepackage{amsmath}
\usepackage{amssymb}
\usepackage{geometry}
\geometry{margin=1in}
\usepackage{hyperref}
\hypersetup{colorlinks=true, linkcolor=blue, citecolor=blue, urlcolor=blue}

\title{Non-Equilibrium Atmospheric Chemistry Plan}
\author{LifeOrig Project}
\date{15 May 2026}

\begin{document}

\maketitle

\section{Introduction}

This document summarizes the current plan for non-equilibrium atmospheric chemistry in the LifeOrig project. The goal is to define the chemical inventory and the key processes to add after the equilibrium atmospheric configuration and radiative-convective structure are established.

The planned non-equilibrium chemistry sources are:
\begin{enumerate}
  \item photochemistry,
  \item geochemical flux and volcanic degassing,
  \item gravity-driven atmospheric escape,
  \item exogenous delivery (meteorites / impacts),
  \item surface/ocean-atmosphere exchange,
  \item electrical and impact shock chemistry,
  \item atmospheric transport, mixing, and aerosol/cloud chemistry.
\end{enumerate}

\section{Initial Chemical Inventory}

The initial equilibrium composition for early Earth is best represented by a background dominated by:
\begin{itemize}
  \item $\mathrm{N_2}$,
  \item $\mathrm{H_2O}$,
  \item $\mathrm{CO_2}$,
  \item $\mathrm{CH_4}$.
\end{itemize}

This set contains the major reservoirs of C, H, O, and N and defines the bulk atmospheric stock available for prebiotic evolution. Trace species such as $\mathrm{HCN}$, $\mathrm{NH_3}$, $\mathrm{H_2S}$, $\mathrm{CO}$, and nitrogen oxides can be generated by non-equilibrium processing.

\section{Non-Equilibrium Process Framework}

The general species continuity equation for an atmospheric species $i$ is:
\begin{equation}
\frac{\partial n_i}{\partial t} + \nabla \cdot \mathbf{\Phi}_i = P_i - L_i + S_i,
\label{eq:continuity}
\end{equation}

where:
\begin{itemize}
  \item $n_i$ is the number density of species $i$,
  \item $\mathbf{\Phi}_i$ is the transport flux,
  \item $P_i$ is the local production rate,
  \item $L_i$ is the loss rate,
  \item $S_i$ is any external source term.
\end{itemize}

In a one-dimensional vertical model with altitude $z$, this becomes:
\begin{equation}
\frac{\partial n_i}{\partial t} + \frac{\partial \Phi_i}{\partial z} = P_i - L_i + S_i.\label{eq:1d}
\end{equation}

\subsection{Photochemistry}

Photochemical production and loss are normally expressed as a sum of chemical reactions:
\begin{equation}
P_i = \sum_r k_r \prod_j n_j, \qquad L_i = \sum_s k_s \prod_j n_j,
\end{equation}

where $k_r$ are reaction rate coefficients and the products run over reactant species. For photolysis reactions, the rate coefficient is:
\begin{equation}
J = \int \sigma(\lambda) \, \phi(\lambda) \, F(\lambda) \, d\lambda,
\end{equation}

with $\sigma(\lambda)$ the absorption cross section, $\phi(\lambda)$ the quantum yield, and $F(\lambda)$ the actinic flux.

Important photochemical pathways for early Earth include:
\begin{itemize}
  \item $\mathrm{CO_2 + h\nu \rightarrow CO + O}$,
  \item $\mathrm{CH_4 + h\nu \rightarrow CH_3 + H}$,
  \item $\mathrm{N_2}$-bearing radical chemistry leading to $\mathrm{HCN}$ and $\mathrm{NH_3}$.
\end{itemize}

\subsection{Geochemical Flux and Volcanic Degassing}

Geochemical sources deliver reduced and oxidized species to the atmosphere at the surface boundary. A source flux can be represented as:
\begin{equation}
S_i^{\mathrm{geo}} = F_i^{\mathrm{deg}} + F_i^{\mathrm{hydro}},
\end{equation}

where $F_i^{\mathrm{deg}}$ is volcanic/tectonic outgassing and $F_i^{\mathrm{hydro}}$ is hydrothermal flux. These fluxes can be computed from rock-water equilibrium and kinetic modeling using a PHREEQC-style geochemistry module.

\subsection{Gravity-Driven Atmospheric Escape}

Atmospheric escape removes species from the top of the atmosphere. For thermal (Jeans) escape:
\begin{equation}
\Phi_i^{\mathrm{esc}} = n_i v_{\mathrm{th}} \left(1 + \frac{\lambda_i}{2}\right) e^{-\lambda_i},
\end{equation}

with $v_{\mathrm{th}} = \sqrt{\frac{2 k_B T}{m_i}}$ and $\lambda_i = \frac{G M m_i}{k_B T r}$ the escape parameter. Nonthermal escape processes such as sputtering and photochemical escape may also contribute.

\subsection{Exogenous Delivery}

Delivery of extraterrestrial material provides an external source:
\begin{equation}
S_i^{\mathrm{exo}} = F^{\mathrm{imp}} \, X_i^{\mathrm{imp}},
\end{equation}

where $F^{\mathrm{imp}}$ is the mass flux of impactors and $X_i^{\mathrm{imp}}$ is the mass fraction of species $i$ in the delivered material.

\subsection{Surface and Ocean Exchange}

Surface and ocean reservoirs can scavenge or release atmospheric species. A simple wet/dry deposition flux is:
\begin{equation}
F_i^{\mathrm{surf}} = -k_i^{\mathrm{dep}} n_i + k_i^{\mathrm{rev}} n_{\mathrm{surf},i},
\end{equation}

with deposition coefficient $k_i^{\mathrm{dep}}$ and surface reservoir coupling $k_i^{\mathrm{rev}}$.

\subsection{Electrical and Impact Shock Chemistry}

High-energy processes produce radicals and fixed nitrogen. A source term can be parameterized as:
\begin{equation}
P_i^{\mathrm{elec}} = \alpha_i I_{\mathrm{elec}} + \beta_i I_{\mathrm{imp}},
\end{equation}

where $I_{\mathrm{elec}}$ is lightning energy deposition and $I_{\mathrm{imp}}$ is impact shock intensity.

\subsection{Transport and Aerosols}

Vertical transport combines eddy diffusion and advection:
\begin{equation}
\Phi_i = -K_{zz} \frac{\partial n_i}{\partial z} + v_z n_i.
\end{equation}

Aerosol and cloud chemistry introduce heterogeneous reactions and can modify photolysis rates and deposition.

\section{Radiative-Convective and Equilibrium Dynamics}

The non-equilibrium chemistry should be built on a stable atmospheric state with a radiative-convective temperature profile. The atmospheric thermal structure can be solved from:
\begin{equation}
\frac{\partial T}{\partial z} = \left( \frac{\partial T}{\partial z} \right)_{\mathrm{rad}} \text{ or } \left( \frac{\partial T}{\partial z} \right)_{\mathrm{ad}},
\end{equation}

with radiative equilibrium where heating and cooling balance and convective adjustments where the temperature lapse rate is superadiabatic.

\section{Suggested Workflow}

\begin{enumerate}
  \item Compute the equilibrium atmospheric configuration and temperature profile.
  \item Implement the radiative-convective solver to produce a physically consistent $T(z)$ profile.
  \item Add photochemical source/loss terms based on the major inventory.
  \item Add geochemical surface fluxes from volcanic and hydrothermal sources.
  \item Add escape and exogenous delivery terms.
  \item Add surface/ocean exchange, lightning/impact chemistry, and transport/aerosol processes.
\end{enumerate}

\section{Notes}

The initial major inventory $\{\mathrm{N_2},\,\mathrm{H_2O},\,\mathrm{CO_2},\,\mathrm{CH_4}\}$ is a reasonable bulk starting point, but the full prebiotic chemical stock requires secondary trace species formed by non-equilibrium chemistry. This document is intentionally high-level and will be refined as the atmospheric dynamics and radiative-convective solver are completed.

\end{document}
